\section{Discussion}
\label{sec:discussion}

\subsection{What the agreement does and does not establish}

Every implemented component agrees with its independent reference to the accuracy of the reference: the five ECB statistics to \ecbAllMax{}, the large-portfolio law to \lhpMax{}, the CreditRisk+ recursion to \crpIndMax{} against an independent recursion, the copula sampler to within a $|z|$ below one against an independent simulator, and the interest rate present values to \irrQlMax{} against QuantLib. These are necessary conditions. A defect that the engine and its reference share, a wrong reading of a specification, or an error in a region of the input space that no reference visits, would pass every check of Section~\ref{sec:results}. The checks are strongest where the reference differs most from the component: the transcription of the ECB statistics uses loops and a quadrature where the engine uses vectorised code and a library distribution function, and an independent simulator draws its own random numbers. They are weakest where the reference is the same arithmetic, as in the IRB formula and the QuantLib present value, and the manuscript says so where it reports them. The measured scope of the catalogue is narrower than a reading of its title suggests. The engine implements \ecbImpl{} of the \ecbSections{} statistics that the ECB instructions fix, all of them on the probability of default, and none on the loss given default, the credit conversion factor or the exposure at default beyond a related two-sided paired $t$-test; a validator who needs the ECB's loss-side statistics as specified does not find them in this engine.

\subsection{The Jeffreys test as an assembly, not a new statistic}

The Jeffreys test is a Beta distribution function evaluated at the applied probability, and Proposition~\ref{prop:sandwich} places it between two binomial tails. On the grades of the credit-card holdout this means the generic tests give nearly the same decisions: \flagJeffFive{} cells flagged against \flagBinFive{} at 5\,\% and none different at 1\,\%, and the statsmodels interval reproduces the engine's decision in every cell. The reason to use the engine's test is therefore not the arithmetic. It is that the grade-by-grade tests and the portfolio test are produced together, with the obligor-weighted portfolio probability that the instructions require, in the result record that carries the null hypothesis, the primary source and the policy, so that a report assembled from them is the report the instructions describe. A reader who expects a different statistic from the Jeffreys test than from the binomial test in a large grade will not find one, and Proposition~\ref{prop:sandwich} says why: the difference is bounded by the probability of the observed count. The substitution is not free in a small grade, where the largest difference between the two $p$-values is \maxPGap{}; but it is conservative, because the binomial test cannot flag a grade that the ECB test passes.

\subsection{Why the five layers are different layers}

The ablation of Section~\ref{sec:ablation} shows that each layer sees a different part of the engine. The analytic anchors see the distribution functions and the recursion and are blind to the sampler and to the tests. The published vectors and the ECB transcription see constants, priors and coefficients that a formula carries, and nothing that depends on a seed. The regeneration of the committed rung sees every change of the sampler's sample path, including those that are not defects, and for that reason it is a regression check and not a correctness check; it was the only layer to find the loading defect, whose size is given in Section~\ref{sec:ablation}. The known-truth simulator sees a biased estimator that no arithmetic check can see, since the defect is in the estimator's output and not in a formula the other layers evaluate. The cross-implementation layer sees the compounding rule that a formula-level check does not exercise, and it has a resolution that Proposition~\ref{prop:power} quantifies. No layer is redundant in the \ablReal{} defects tried, and no layer is sufficient: \ablNone{} real defect is found by none, and the others are found by between one and two layers.

\subsection{Negative results and failure cases}

Four results go against a simple reading. The correlation interval does not cover as claimed in the cell at the lowest correlation and the shortest series, with a coverage of \gridMlWorstCov{} against a nominal $0.95$; the moments estimator is outside the band in \gridMomFail{} of \gridMomCells{} cells, so its interval should not be read as a 95\,\% interval anywhere in the grid. The engine's absent ECB statistics are a negative result about the catalogue, not about any test. The Student $t$ importance sampler does not exist and the rung says so. And the one comparison in which a defect of $6\,\%$ in the correlation passes a cross-implementation check at the size run in the ablation is a statement about that layer's resolution and about the choice of its size, not about the engine. A validator who uses the layers should read them as the table of Section~\ref{sec:ablation} suggests: the five together find what each alone misses, and the absence of a finding is a statement about the references used.
